% GENERATED FILE. Edit canonical lessons, then run npm run generate:books.
% canonical-source-hash: fnv1a64:dbcd617b35173839
% canonical-lessons: BN-C144-lajuk, BN-C144-mishuk, BN-C144-chupchap, BN-C144-molayem, BN-C144-khoskhose, BN-R144-first-pass-describing-words, BN-R144-second-pass-describing-people

\chapter{Shy, Sociable, Quiet, Smooth, Rough}
\label{ch:bn-shy-sociable-quiet-smooth-rough}
\hlchaptermodality{144}

\begin{chapteropening}
\textbf{By the end of this chapter:} \emph{I can say shy, sociable, quiet, smooth and rough.}

Complete the last lesson of chapter 144: I can say shy, sociable, quiet, smooth and rough.
\end{chapteropening}

\section[\texorpdfstring{\bn{লাজুক} (lājuk)}{লাজুক (lājuk)}]{\bn{লাজুক} (lājuk) — shy}

\label{lesson:BN-C144-lajuk}



\begin{quote}
Before the new one: say the Bengali for small, then the Bengali for cheerful.
\end{quote}

\subsection*{You'll want to know: \bn{লাজুক}}
\textbf{\bn{লাজুক}} — \emph{lājuk} — \textquotedblleft{}shy\textquotedblright{}.

\begin{grammarlens}[title={What you've built}]
One more word for this chapter.
\end{grammarlens}

\subsection*{Your turn}
\begin{itemize}
\raggedright
  \item \emph{Say it:} \emph{lājuk}
  \item \emph{Say it:} \emph{lājuk}, once more
  \item \emph{Say it:} say \emph{lājuk}
  \item \emph{Recall:} say the Bengali for small, then the Bengali for cheerful, then say \emph{lājuk} again
\end{itemize}

\begin{tcolorbox}[breakable,colback=teal!4,colframe=teal!35!black,title={Before you move on}]
What does \bn{লাজুক} mean? (\textquotedblleft{}Shy\textquotedblright{}.) Say it once more.
\end{tcolorbox}
\section[\texorpdfstring{\bn{মিশুক} (mishuk)}{মিশুক (mishuk)}]{\bn{মিশুক} (mishuk) — sociable}

\label{lesson:BN-C144-mishuk}



\begin{quote}
Before the new one: say the Bengali for cheerful, then the Bengali for shy.
\end{quote}

\subsection*{You'll want to know: \bn{মিশুক}}
\textbf{\bn{মিশুক}} — \emph{mishuk} — \textquotedblleft{}sociable\textquotedblright{}.

\begin{grammarlens}[title={What you've built}]
One more word for this chapter.
\end{grammarlens}

\subsection*{Your turn}
\begin{itemize}
\raggedright
  \item \emph{Say it:} \emph{mishuk}
  \item \emph{Say it:} \emph{mishuk}, once more
  \item \emph{Say it:} say \emph{mishuk}
  \item \emph{Recall:} say the Bengali for cheerful, then the Bengali for shy, then say \emph{mishuk} again
\end{itemize}

\begin{tcolorbox}[breakable,colback=teal!4,colframe=teal!35!black,title={Before you move on}]
What does \bn{মিশুক} mean? (\textquotedblleft{}Sociable\textquotedblright{}.) Say it once more.
\end{tcolorbox}
\section[\texorpdfstring{\bn{চুপচাপ} (chupchāp)}{চুপচাপ (chupchāp)}]{\bn{চুপচাপ} (chupchāp) — quiet}

\label{lesson:BN-C144-chupchap}



\begin{quote}
Before the new one: say the Bengali for shy, then the Bengali for sociable.
\end{quote}

\subsection*{You'll want to know: \bn{চুপচাপ}}
\textbf{\bn{চুপচাপ}} — \emph{chupchāp} — \textquotedblleft{}quiet\textquotedblright{}.

\begin{grammarlens}[title={What you've built}]
One more word for this chapter.
\end{grammarlens}

\subsection*{Your turn}
\begin{itemize}
\raggedright
  \item \emph{Say it:} \emph{chupchāp}
  \item \emph{Say it:} \emph{chupchāp}, once more
  \item \emph{Say it:} say \emph{chupchāp}
  \item \emph{Recall:} say the Bengali for shy, then the Bengali for sociable, then say \emph{chupchāp} again
\end{itemize}

\begin{tcolorbox}[breakable,colback=teal!4,colframe=teal!35!black,title={Before you move on}]
What does \bn{চুপচাপ} mean? (\textquotedblleft{}Quiet\textquotedblright{}.) Say it once more.
\end{tcolorbox}
\section[\texorpdfstring{\bn{মোলায়েম} (molāyem)}{মোলায়েম (molāyem)}]{\bn{মোলায়েম} (molāyem) — smooth}

\label{lesson:BN-C144-molayem}



\begin{quote}
Before the new one: say the Bengali for sociable, then the Bengali for quiet.
\end{quote}

\subsection*{You'll want to know: \bn{মোলায়েম}}
\textbf{\bn{মোলায়েম}} — \emph{molāyem} — \textquotedblleft{}smooth\textquotedblright{}.

\begin{grammarlens}[title={What you've built}]
One more word for this chapter.
\end{grammarlens}

\subsection*{Your turn}
\begin{itemize}
\raggedright
  \item \emph{Say it:} \emph{molāyem}
  \item \emph{Say it:} \emph{molāyem}, once more
  \item \emph{Say it:} say \emph{molāyem}
  \item \emph{Recall:} say the Bengali for sociable, then the Bengali for quiet, then say \emph{molāyem} again
\end{itemize}

\begin{tcolorbox}[breakable,colback=teal!4,colframe=teal!35!black,title={Before you move on}]
What does \bn{মোলায়েম} mean? (\textquotedblleft{}Smooth\textquotedblright{}.) Say it once more.
\end{tcolorbox}
\section[\texorpdfstring{\bn{খসখসে} (khôskhôse)}{খসখসে (khôskhôse)}]{\bn{খসখসে} (khôskhôse) — rough}

\label{lesson:BN-C144-khoskhose}



\begin{quote}
Before the new one: say the Bengali for quiet, then the Bengali for smooth.
\end{quote}

\subsection*{You'll want to know: \bn{খসখসে}}
\textbf{\bn{খসখসে}} — \emph{khôskhôse} — \textquotedblleft{}rough\textquotedblright{}.

\begin{grammarlens}[title={What you've built}]
One more word for this chapter.
\end{grammarlens}

\subsection*{Your turn}
\begin{itemize}
\raggedright
  \item \emph{Say it:} \emph{khôskhôse}
  \item \emph{Say it:} \emph{khôskhôse}, once more
  \item \emph{Say it:} say \emph{khôskhôse}
  \item \emph{Recall:} say the Bengali for quiet, then the Bengali for smooth, then say \emph{khôskhôse} again
\end{itemize}

\begin{tcolorbox}[breakable,colback=teal!4,colframe=teal!35!black,title={Before you move on}]
What does \bn{খসখসে} mean? (\textquotedblleft{}Rough\textquotedblright{}.) Say it once more.
\end{tcolorbox}
\section[(dialogue)]{Describing words — a first pass}

\label{lesson:BN-R144-first-pass-describing-words}



\begin{quote}
Say the words for smooth and rough.
\end{quote}

\subsection*{Your turn}
Say these aloud:

\begin{itemize}
\raggedright
  \item busy, clever, foolish, kind, brave
  \item timid, calm, urgent, straight, crooked
  \item sharp, blunt, fresh, stale, raw
\end{itemize}

\begin{tcolorbox}[breakable,colback=teal!4,colframe=teal!35!black,title={Before you move on}]
Busy? (\textbf{\bn{ব্যস্ত}}, \emph{byāstô}.) Clever? (\textbf{\bn{চালাক}}, \emph{chālāk}.) Foolish? (\textbf{\bn{বোকা}}, \emph{bokā}.) Kind? (\textbf{\bn{দয়ালু}}, \emph{dôyālu}.) Brave? (\textbf{\bn{সাহসী}}, \emph{sāhôsi}.) Timid? (\textbf{\bn{ভীতু}}, \emph{bhītu}.) Calm? (\textbf{\bn{শান্ত}}, \emph{shāntô}.) Urgent? (\textbf{\bn{জরুরি}}, \emph{jôruri}.) Straight? (\textbf{\bn{সোজা}}, \emph{sojā}.) Crooked? (\textbf{\bn{বাঁকা}}, \emph{bā̃kā}.) Sharp? (\textbf{\bn{ধারালো}}, \emph{dhārālo}.) Blunt? (\textbf{\bn{ভোঁতা}}, \emph{bhõtā}.) Fresh? (\textbf{\bn{টাটকা}}, \emph{ṭāṭkā}.) Stale? (\textbf{\bn{বাসি}}, \emph{bāsi}.) Raw? (\textbf{\bn{কাঁচা}}, \emph{kā̃chā}.)
\end{tcolorbox}
\section[(dialogue)]{Describing people — a second pass}

\label{lesson:BN-R144-second-pass-describing-people}



\begin{quote}
Say the words for ripe and funny.
\end{quote}

\subsection*{Your turn}
Say these aloud:

\begin{itemize}
\raggedright
  \item ripe, funny, comfortable, dangerous, safe
  \item dear, famous, clean, sturdy, weak
  \item powerful, tasty, fairly big, small, cheerful
  \item shy, sociable, quiet, smooth, rough
\end{itemize}

\begin{tcolorbox}[breakable,colback=teal!4,colframe=teal!35!black,title={Before you move on}]
Ripe? (\textbf{\bn{পাকা}}, \emph{pākā}.) Funny? (\textbf{\bn{মজার}}, \emph{môjār}.) Comfortable? (\textbf{\bn{আরামদায়ক}}, \emph{ārāmdāyôk}.) Dangerous? (\textbf{\bn{বিপজ্জনক}}, \emph{bipôjjônôk}.) Safe? (\textbf{\bn{নিরাপদ}}, \emph{nirāpôd}.) Dear? (\textbf{\bn{প্রিয়}}, \emph{priyô}.) Famous? (\textbf{\bn{বিখ্যাত}}, \emph{bikhyātô}.) Clean? (\textbf{\bn{পরিচ্ছন্ন}}, \emph{pôrichchhônnô}.) Sturdy? (\textbf{\bn{মজবুত}}, \emph{môjbut}.) Weak? (\textbf{\bn{দুর্বল}}, \emph{durbôl}.) Powerful? (\textbf{\bn{শক্তিশালী}}, \emph{shôktishāli}.) Tasty? (\textbf{\bn{সুস্বাদু}}, \emph{susbādu}.) Fairly big? (\textbf{\bn{বড়সড়}}, \emph{bôṛôsôṛô}.) Small? (\textbf{\bn{ছোটখাটো}}, \emph{chhoṭôkhāṭo}.) Cheerful? (\textbf{\bn{হাসিখুশি}}, \emph{hāsikhushi}.) Shy? (\textbf{\bn{লাজুক}}, \emph{lājuk}.) Sociable? (\textbf{\bn{মিশুক}}, \emph{mishuk}.) Quiet? (\textbf{\bn{চুপচাপ}}, \emph{chupchāp}.) Smooth? (\textbf{\bn{মোলায়েম}}, \emph{molāyem}.) Rough? (\textbf{\bn{খসখসে}}, \emph{khôskhôse}.)
\end{tcolorbox}
